\section{Methodology}
\label{sec:method}

In this section, we present the mathematical and architectural formulation of our proposed \textbf{Q-PolyMerge} framework. We detail our polynomial subspace parameterization, our symmetric uniform post-training quantization operator, our test-time adaptation objective, and our two optimization pathways.

\subsection{Preliminaries \& Task Vectors}
Let $\Theta_{\text{base}} \in \mathbb{R}^M$ denote the parameters of a pre-trained base backbone model (e.g., a Vision Transformer). Let $\{\Theta_k\}_{k=1}^K$ represent a set of $K$ independent task-specific expert checkpoints, each fine-tuned from the identical base checkpoint $\Theta_{\text{base}}$ on a distinct downstream dataset. We partition the backbone parameters into $L$ sequential layer-specific parameter blocks (such as attention weight matrices, feed-forward layers, or projection matrices). For our Vision Transformer (ViT-Tiny) setup, the backbone is partitioned into $L=14$ blocks, where Layer 0 represents the patch embedding layer, Layers 1--12 correspond to the 12 sequential Transformer blocks, and Layer 13 represents the final layer normalization block. 

For each layer block $l \in \{0, 1, \dots, L-1\}$, the expert task vector $\mathbf{\Delta}_{k, l} \in \mathbb{R}^{M_l}$ is defined as the parameter difference between the expert weights and the base weights:
\begin{equation}
\mathbf{\Delta}_{k, l} = \Theta_{k, l} - \Theta_{\text{base}, l}
\end{equation}
The classic task arithmetic framework \cite{ilharco2022editing} merges these task vectors into a single unified multi-task model using static, uniform scaling coefficients. Under test-time adaptation, we seek layer-wise and expert-wise coefficients $\lambda_{k, l} \in \mathbb{R}$ to form a continuous merged model before quantization:
\begin{equation}
\label{eq:continuous_merge}
\Theta_{\text{merged}, l} = \Theta_{\text{base}, l} + \sum_{k=1}^K \lambda_{k, l} \mathbf{\Delta}_{k, l}
\end{equation}

\subsection{Continuous Polynomial Subspace Parameterization}
Under test-time adaptation, optimizing individual coefficients $\lambda_{k, l}$ independently for each layer $l$ and task $k$ requires learning $L \times K$ parameters. On small, unlabeled on-device calibration batches, this high dimensionality allows the optimizer to overfit to statistical noise in the calibration stream. This \emph{Overfitting-Optimizer Paradox} leads to highly jagged layer-wise trajectories that generalize poorly to held-out test data.

To resolve this, Q-PolyMerge constrains the merging coefficients to lie within a continuous, low-degree polynomial subspace of normalized layer depth. We define the layer coefficient $\lambda_{k, l}(\boldsymbol{\alpha})$ as a polynomial of degree $d$ evaluated at the normalized depth $\bar{l} = \frac{l}{L-1} \in [0, 1]$:
\begin{equation}
\label{eq:polynomial_mapping}
\lambda_{k, l}(\boldsymbol{\alpha}) = \sum_{j=0}^d \alpha_{k, j} \cdot \left( \frac{l}{L-1} \right)^j
\end{equation}
where $\boldsymbol{\alpha} = \{ \alpha_{k, j} \} \in \mathbb{R}^{K \times (d+1)}$ represents the small set of learnable continuous polynomial parameters. Normalizing the layer indices to the compact domain $[0, 1]$ is a crucial operational detail: it guarantees numerical stability by preventing polynomial terms from diverging in deep architectures, and it ensures scale invariance across networks of varying depths.

By substituting Equation \ref{eq:polynomial_mapping} into Equation \ref{eq:continuous_merge}, the continuous merged weights at layer $l$ are governed by:
\begin{equation}
\Theta_{\text{merged}, l}(\boldsymbol{\alpha}) = \Theta_{\text{base}, l} + \sum_{k=1}^K \left[ \sum_{j=0}^d \alpha_{k, j} \left( \frac{l}{L-1} \right)^j \right] \mathbf{\Delta}_{k, l}
\end{equation}
This projects the learnable dimensions from $L \times K$ down to just $(d+1) \times K$. For a standard network with $L=14$ layers and $K=4$ tasks using a quadratic polynomial ($d=2$), Q-PolyMerge reduces the optimization space from 56 parameters to just 12, a \textbf{78.6\% search space reduction}. Bounding the coefficients to a smooth polynomial acts as a mathematical low-pass filter, filtering out high-frequency optimization noise and preventing overfitting.

While the standard monomial basis $\{1, \bar{l}, \bar{l}^2, \dots\}$ is numerically stable and highly performant for low-degree polynomial constraints ($d \le 3$), scaling to extremely deep networks (such as 32-to-80 layer foundation models) can cause monomial bases to become ill-conditioned. In such deep regimes, alternative orthogonal polynomial bases—such as Chebyshev polynomials of the first kind $T_j(\cdot)$—can be seamlessly integrated to prevent Runge's boundary oscillations. We detail these orthogonal scaling pathways and formulations in Appendix~\ref{sec:appendix_chebyshev_scaling}.

\subsection{Symmetric Uniform Post-Training Quantization}
We compress the merged model to low-bit formats using a symmetric uniform post-training quantization operator $q(\cdot)$.
For standard 8-bit quantization (INT8), we employ per-tensor uniform quantization. For aggressive 4-bit quantization (INT4), to protect coordinate alignment and mode connectivity from extreme outlier values, we mandate per-channel uniform weight quantization.

The per-channel quantization operator is formulated as:
\begin{equation}
q(W_{c, \cdot}) = \text{clamp}\left( \text{round}\left( \frac{W_{c, \cdot}}{S_c} \right), -2^{b-1}, 2^{b-1}-1 \right) \cdot S_c
\end{equation}
where $W_{c, \cdot}$ represents the weights of output channel $c$, $b \in \{4, 8\}$ is the bit-width, and $S_c$ is the per-channel scale factor calculated from the absolute maximum weight value in that channel:
\begin{equation}
S_c = \frac{\max_i(|W_{c, i}|)}{2^{b-1} - 1}
\end{equation}
The quantized merged weights of the model are thus represented as:
\begin{equation}
\Theta^q_{\text{merged}, l}(\boldsymbol{\alpha}) = q\left( \Theta_{\text{merged}, l}(\boldsymbol{\alpha}) \right)
\end{equation}

\subsection{Test-Time Optimization Objective}
During test-time adaptation, the network receives a streaming calibration batch of unlabeled target images. Adhering to the principles of test-time entropy minimization, we optimize the continuous polynomial parameters $\boldsymbol{\alpha}$ by minimizing the average Shannon entropy of the model's predictions:
\begin{equation}
\mathcal{L}_{\text{TTA}}(\boldsymbol{\alpha}) = \sum_{k=1}^K \mathbb{E}_{x \sim \mathcal{D}_k^{\text{unlabeled}}} \left[ H\left( f_{\Theta^q_{\text{merged}}(\boldsymbol{\alpha})}(x) \right) \right]
\end{equation}
where $f_{\Theta^q_{\text{merged}}(\boldsymbol{\alpha})}(x)$ is the logit distribution predicted by the fully quantized network, and $H(\mathbf{p}) = -\sum_{c=1}^C p_c \log (p_c + \epsilon)$ represents the Shannon entropy (with a stability constant $\epsilon = 10^{-8}$). Minimizing entropy forces the quantized model's representations to become highly confident, aligning task distributions and mitigating the representational misalignment caused by quantization noise.

To prevent the continuous polynomial parameters $\boldsymbol{\alpha}$ from learning arbitrarily large scaling trajectories that could scale merged weights far outside normal boundaries and trigger severe clipping under the physical quantization operator, we can incorporate a trajectory clamping function or L2 regularization. Specifically, we can restrict the learned layer coefficients $\lambda_{k, l}(\boldsymbol{\alpha})$ to a stable, physically meaningful interpolation interval $[-\gamma, 1+\gamma]$ (where typically $\gamma = 0.5$):
\begin{equation}
\tilde{\lambda}_{k, l}(\boldsymbol{\alpha}) = \max\left(-\gamma, \min\left(1+\gamma, \lambda_{k, l}(\boldsymbol{\alpha})\right)\right)
\end{equation}
Alternatively, an L2 weight decay term $\frac{\beta}{2} \|\boldsymbol{\alpha}\|_F^2$ can be added directly to the test-time loss $\mathcal{L}_{\text{TTA}}$ to penalize extreme coefficient scales. In our physical PyTorch implementation, we implement both constraints: we clamp the layer coefficients $\lambda_{k, l}$ to the stable interval $[-0.5, 1.5]$ (using $\gamma = 0.5$) and incorporate L2 weight decay on $\boldsymbol{\alpha}$ (with $\beta = 10^{-4}$). Bounding the coefficient trajectories in this manner serves as an essential stability safeguard, protecting the dynamic range of weight values prior to uniform PTQ rounding and further stabilizing adaptation.

\subsection{First-Order Optimization via Straight-Through Estimator}
Because the $\text{round}$ function in the quantization operator has a derivative of zero almost everywhere, the standard backpropagation of gradients is blocked. We bypass this using the Straight-Through Estimator (STE) \cite{bengio2013estimating}, which approximates the derivative of the rounding operator as the identity mapping during the backward pass:
\begin{equation}
\frac{\partial q(W)}{\partial W} \approx \mathbf{I}
\end{equation}
Applying the multi-variable chain rule, the gradient of the test-time loss with respect to our learnable polynomial parameter $\alpha_{k, j}$ is computed as:
\begin{equation}
\frac{\partial \mathcal{L}_{\text{TTA}}}{\partial \alpha_{k, j}} = \sum_{l=0}^{L-1} \frac{\partial \mathcal{L}_{\text{TTA}}}{\partial \Theta^q_{\text{merged}, l}} \cdot \frac{\partial \Theta^q_{\text{merged}, l}}{\partial \lambda_{k, l}} \cdot \frac{\partial \lambda_{k, l}}{\partial \alpha_{k, j}}
\end{equation}
Under the STE approximation, the gradient flows directly through the quantization operator to scale the task vectors:
\begin{equation}
\frac{\partial \mathcal{L}_{\text{TTA}}}{\partial \alpha_{k, j}} \approx \sum_{l=0}^{L-1} \left\langle \frac{\partial \mathcal{L}_{\text{TTA}}}{\partial \Theta^q_{\text{merged}, l}}, \mathbf{\Delta}_{k, l} \right\rangle \cdot \left( \frac{l}{L-1} \right)^j
\end{equation}
The continuous parameters $\boldsymbol{\alpha}$ are iteratively updated using the first-order Adam optimizer on the device.

\subsection{Zero-Order Optimization via 1+1 Evolution Strategy}
In real-world microcontroller deployments, backpropagation is highly impractical because caching activations and storing gradient graphs consumes massive amounts of volatile memory (SRAM), which is severely limited on the edge. 

To overcome this constraint, we implement a derivative-free, zero-order optimization pathway based on the 1+1 Evolution Strategy (1+1 ES). This treats the quantized model as a black-box oracle. At each iteration, we perturb the current parameters:
\begin{equation}
\boldsymbol{\alpha}_{\text{candidate}} = \boldsymbol{\alpha}_{\text{active}} + \sigma \cdot \mathbf{z}
\end{equation}
where $\mathbf{z} \sim \mathcal{N}(\mathbf{0}, \mathbf{I})$ and $\sigma$ represents an adaptive step size. We evaluate the candidate model's entropy on the calibration stream. If $\mathcal{L}_{\text{TTA}}(\boldsymbol{\alpha}_{\text{candidate}}) < \mathcal{L}_{\text{TTA}}(\boldsymbol{\alpha}_{\text{active}})$, we accept the candidate: $\boldsymbol{\alpha}_{\text{active}} \leftarrow \boldsymbol{\alpha}_{\text{candidate}}$, and increase $\sigma$ to accelerate search. Otherwise, we reject it and shrink $\sigma$.

Under unconstrained search (Q-Merge 1+1 ES), optimizing in high-dimensional space ($L \times K$ parameters) is extremely inefficient due to the curse of dimensionality, leading to search failure. However, by projecting the search space down to our continuous polynomial parameters (just $(d+1) \times K$ dimensions), the search becomes highly efficient, enabling zero-order optimization to match or exceed gradient methods. Since 1+1 ES only requires standard inference forward passes, it incurs \textbf{zero activation caching, zero gradient overhead, and zero backward-pass memory footprints}, making adaptive on-device low-bit merging physically viable for the first time.

\subsection{Integer-Weight Edge Pipeline}
To deliver on strict storage and memory-bandwidth constraints where weight transfer dominates energy consumption, we post-hoc quantize the task-specific linear classification heads to 8-bit precision (INT8) after completing the test-time optimization of $\boldsymbol{\alpha}$. This yields a \textbf{fully integer-quantized weight pipeline} (W8 or W4 weight representation across both backbone and classification heads), completely removing the need to store any floating-point weight parameters on disk or flash. 

While intermediate activations, layer normalizations, and attention softmax calculations are maintained in FP16/FP32 formats during inference to preserve representational precision (forming a hybrid-precision pipeline), storing 100\% of the parameters as low-bit integers dramatically cuts DRAM-to-SRAM communication bottlenecks, which are the primary drivers of on-device latency and energy dissipation. 

We acknowledge that in extremely low-power and low-cost microcontroller environments lacking dedicated hardware Floating-Point Units (FPUs) or vector engines supporting FP16, executing activations in floating-point format requires software emulation, which incurs substantial latency and energy penalties. A truly edge-optimal execution pipeline would require transitioning to a \textbf{fully-integerized arithmetic pipeline} (e.g., W8A8 or W4A8 execution), where intermediate activations are also quantized post-hoc, and floating-point operations (such as layer normalization and attention softmax) are replaced with integer shift-and-add arithmetic (using libraries like CMSIS-NN or IREE). 

While implementing fully-integerized activation operators is beyond our current scope, our proposed Q-PolyMerge framework is fundamentally orthogonal to the precision of the activation execution: the continuous polynomial constraint and test-time optimization formulation apply identically under W8A8 or W4A8 integer math, making Q-PolyMerge a highly valuable and adaptable prior for next-generation, integer-only hardware execution.
